\section{A duty-cycle model of the detection pipeline}
\label{sec:sysmodel}

A deliberately simple model organizes the measurements. Its three
parts describe how engine time depends on the GPU clock, how each
pipeline schedule sets the GPU's duty cycle, and how the governor picks
a clock for that duty cycle. A fourth relation splits energy into an
idle floor and a marginal part. \Cref{fig:framework}(b) summarizes the loop. Each part is checked
against measurements in \cref{sec:diagnosis,sec:ladder,sec:stream}.

\textbf{Pipeline stages.}
For each image, the CPU runs a decode-and-letterbox stage $D$ and,
for float-input engines only, a host normalization $N$. The GPU runs
the engine $G$, including copies, which take under 1\,ms on this
unified-memory device. A row construction stage $R$ converts the $300$
output rows to image coordinates. Let $M$ denote the remaining
per-image work of the main thread (buffer fill, enqueue, synchronization).

\textbf{Engine time.}
For a fixed engine, the GPU time at clock $f$ is modeled as
\begin{equation}
  G(f) = G_{\max}\,\frac{f_{\max}}{f},
  \label{eq:gf}
\end{equation}
where $G_{\max}$ is the time at the maximum clock $f_{\max}$. This
assumes the engine is limited by the GPU clock and not by memory. For
variable-clock runs, the time-weighted harmonic mean of the sampled
clock is used as an approximate summary consistent with \cref{eq:gf}.

\textbf{Schedules and duty cycle.}
The period $T$ between completed images, and thus throughput
$X = 1/T$, depends on which stages lie on the critical path. For
$k$ decode workers:
\begin{align}
  T_{\text{seq}} &= D + N + G(f) + R + M, \label{eq:seq}\\
  T_{\text{prefetch}} &\approx \max\!\big(\tfrac{D+N}{k},\; G(f) + R + M\big), \label{eq:pre}\\
  T_{\text{overlap}} &\approx \max\!\big(\tfrac{D+N}{k},\; R + M,\; G(f)\big). \label{eq:ovl}
\end{align}
The GPU busy fraction is $u = G(f)/T$. In the sequential schedule every
host millisecond is GPU idle time. \Cref{eq:pre} describes the prefetch
transition: host stages leave the critical path, and the
period becomes the maximum of the worker rate and the GPU-serial part.
The stage times in it are clock-dependent, so it serves as a structure
rather than a fitted predictor. \Cref{sec:diagnosis} measures them at a
down-clocked host; at the higher clocks of the prefetch rung, the
decode rate is correspondingly higher. Overlap also removes $R$
(\cref{eq:ovl}).

\textbf{Governor.}
The GPU governor (\texttt{nvhost\_podgov}) ships in NVIDIA's public
Jetson Linux kernel sources~\cite{nvidia2025jetsonkernel}. Its devfreq
sysfs node on the test device exposes only a 25\,ms polling interval
and no load thresholds. The thresholds were not derived from the
source; the governor is instead characterized empirically as a
load-band controller. The driver exposes the load it acts on through
sysfs, so the busy fraction $u$ can be checked directly
(\cref{sec:diagnosis}). The governor raises $f$ while $u$ exceeds an
upper threshold $u_{\text{hi}}$, lowers it while $u$ is below a lower
threshold $u_{\text{lo}}$, and is clamped to $[f_{\min}, f_{\max}]$. A
schedule therefore settles either at an interior clock with
$u \in [u_{\text{lo}}, u_{\text{hi}}]$ or at one of the clamps. The
measurements show a fixed point that this view implies. In
\cref{eq:seq}, $u$ increases as $f$ decreases, because a slower GPU occupies a larger share of the
period. If, even at $f_{\min}$, the sequential pipeline's busy fraction
$G(f_{\min})/T_{\text{seq}}(f_{\min})$ stays below $u_{\text{hi}}$, the
governor has no reason to raise the clock, and the application runs at
the minimum clock indefinitely, even though the GPU is busy
about half of the time (measured loads 0.46--0.66 at the floor,
\cref{sec:diagnosis}).

\textbf{Energy.}
Let $P_{\text{idle}}$ be the module input power with no workload under a
given clock policy, and $P_{\text{board}}$ the mean power during a run.
Energy per image splits as
\begin{equation}
  E = \frac{P_{\text{board}}}{X}
    = \underbrace{\frac{P_{\text{idle}}}{X}}_{\text{idle reference}}
    + \underbrace{\frac{P_{\text{board}} - P_{\text{idle}}}{X}}_{E_{\text{marg}}}.
  \label{eq:energy}
\end{equation}
The first term is an accounting reference: the idle draw of the board
in the measured power state, amortized over throughput. It is not a
universal lower bound, since a different sleep or power-gating regime
could go below it. On a backlog, a faster schedule shrinks the
idle-reference term. For a camera that delivers frames at a fixed rate
$r$ below the pipeline's capacity, the throughput is $X = r$ regardless
of schedule, so
\begin{equation}
  E(r) \approx \frac{P_{\text{idle}}}{r} + E_{\text{marg}}.
  \label{eq:stream}
\end{equation}
The schedule then affects energy only through $E_{\text{marg}}$, so
pipeline optimization should yield large energy savings on a backlog
but only small ones at fixed camera rates. \Cref{sec:stream} tests
this prediction and finds that \cref{eq:stream} with the backlog $E_{\text{marg}}$ substituted is an
empirical upper envelope for the below-capacity cells, not an exact
prediction.
